% Ч1.2 — Среднее арифметическое
% Старая методичка: стр. 216 (старое задание 19)

\begin{tema}{Ч1.2}
\opr{Средним арифметическим} чисел $a_1, a_2, \dots, a_n$ называется число
\[
  m=\frac{a_1+a_2+\dots+a_n}{n}.
\]
В задачах часто удобнее обратная запись: \emph{сумма равна среднему, умноженному на количество},
\[
  a_1+a_2+\dots+a_n=m\cdot n.
\]

\begin{primer}
Средний балл по пяти оценкам равен 4, значит, сумма оценок равна $4\cdot5=20$.
\end{primer}

Связанные темы: математическое ожидание \tref{В3.2} — это «среднее с весами»,
где весами служат вероятности.

\begin{oshibka}
Средняя скорость — \emph{не} среднее арифметическое скоростей:
это весь путь, делённый на всё время (\tref{М1.7}).
\end{oshibka}
\end{tema}
